\documentclass[10pt]{article}
\usepackage[T1]{fontenc}
\usepackage[utf8]{inputenc}
\usepackage[french]{babel}
\usepackage{amsmath,amssymb,mathrsfs,mathpazo}
\usepackage{geometry}
\geometry{a4paper,margin=2cm}
\usepackage{graphicx,array,tabularx,tabularray,colortbl,multirow,multicol,booktabs,enumitem,pifont,ulem,fancybox,diagbox,cancel}
\usepackage{pstricks,pstricks-add,pst-plot,pst-tree,pst-node,pst-text,pst-func,pst-math,pst-bspline,pst-3dplot,pst-eucl,pst-solides3d}
\let\pscalculate\undefined
\usepackage{pst-bezier}
\usepackage{tikz,pgfplots,tkz-tab}
\usetikzlibrary{arrows,arrows.meta,patterns,calc,positioning,trees,decorations.markings,shapes.geometric,shapes.misc,shapes.arrows,intersections,decorations.pathmorphing,decorations.pathreplacing,backgrounds,shadows,babel}
\pgfplotsset{compat=1.18}
\usepgfplotslibrary{fillbetween}
\usepackage[np]{numprint}
\usepackage{listings,makecell,pas-tableur,twoopt}
\let\barre\undefined
\usepackage[pstricks]{tablvar}
\definecolor{darkgreen}{RGB}{0,100,0}
\providecommand{\red}{\color{red}}
\providecommand{\blue}{\color{blue}}
\providecommand{\green}{\color{green}}
\providecommand{\black}{\color{black}}
\providecommand{\white}{\color{white}}
\providecommand{\gray}{\color{gray}}
\providecommand{\cyan}{\color{cyan}}
\providecommand{\magenta}{\color{magenta}}
\providecommand{\yellow}{\color{yellow}}
\providecommand{\violet}{\color{violet}}
\providecommand{\orange}{\color{orange}}
\providecommand{\euro}{\texteuro}
\providecommand{\EUR}{\texteuro}
\providecommand{\ieme}{\textsuperscript{e}}
\providecommand{\degres}{\ensuremath{^\circ}}
\providecommand{\degree}{\ensuremath{^\circ}}
\providecommand{\vv}[1]{\overrightarrow{#1}}
\providecommand{\Coord}[1]{#1}
\providecommand{\decosix}{}
\providecommand{\starredbullet}{\ensuremath{\bullet}}
\pagestyle{empty}
\setlength{\parindent}{0pt}
\usepackage{adjustbox}\DeclareUnicodeCharacter{2236}{:}\DeclareUnicodeCharacter{2190}{\ensuremath{\leftarrow}}
\providecommand{\R}{\mathbb{R}}
\providecommand{\N}{\mathbb{N}}
\providecommand{\Z}{\mathbb{Z}}
\providecommand{\Q}{\mathbb{Q}}
\providecommand{\D}{\mathbb{D}}
\providecommand{\C}{\mathbb{C}}
\providecommand{\e}{{\mathrm{e}}}
\providecommand{\vect}[1]{\overrightarrow{#1}}
\providecommand{\vectt}[1]{\overrightarrow{\mathrm{#1}}}
\providecommand{\vvt}[1]{\overrightarrow{\mathrm{#1}}}
\providecommand{\np}[2][]{{#2~#1}}
\providecommand{\vv}[1]{{\overrightarrow{#1}}}
\renewcommand{\d}{{\mathrm{d}}}
\providecommand{\degres}{{^\circ}}
\providecommand{\degre}{{^\circ}}
\providecommand{\blue}{\color{blue}}
\providecommand{\red}{\color{red}}
\providecommand{\enskip}{\quad}
\providecommand{\llbracket}{\lbrack\!\lbrack}
\providecommand{\rrbracket}{\rbrack\!\rbrack}
\providecommand{\euro}{{\text{€}}}
\providecommand{\no}{{\text{n°}}}
\providecommand{\interval}[2]{{[#1;#2]}}
\providecommand{\intervFO}[2]{{[#1;#2[}}
\providecommand{\dotfill}{\ldots}

\providecommand{\ds}{\displaystyle}
\providecommand{\textding}[1]{#1}
\providecommand{\Oij}{\(\mathrm{O};\vec{\imath},\vec{\jmath}\)}
\providecommand{\Oijk}{\(\mathrm{O};\vec{\imath},\vec{\jmath},\vec{k}\)}
\providecommand{\Ouv}{\(\mathrm{O};\vec{u},\vec{v}\)}
\providecommand{\intervFF}[2]{\ensuremath{[#1;#2]}}
\renewcommand{\euro}{\ifmmode\text{\texteuro}\else\texteuro\fi}
\DeclareUnicodeCharacter{00B0}{\ifmmode^{\circ}\else\textdegree\fi}
\DeclareUnicodeCharacter{20AC}{\ifmmode\text{\texteuro}\else\texteuro\fi}
% Énoncés : Éducation nationale. Transcription et corrigés : APMEP, Denis Vergès et contributeurs. Découpage et classement : CoopMaths, Sébastien Lozano, Éric Elter et contributeurs. Adaptation de la présentation : Maths Explorer.
% CC BY-SA : https://coopmaths.fr/www/about/
% https://coopmaths.fr/alea/static/bac/2025/tex/bac_2025_11_sujet2_caledonie_3_cor.tex
\begin{document}
\noindent\textit{APMEP / CoopMaths. CC BY-SA. Présentation adaptée par Maths Explorer.}\par\medskip
\smallskip

On considère la fonction $f$ définie pour tout réel $x$ par:

\[f(x) = \ln\left ( \text{e}^{\frac{x}{2}}+2 \right ).\]

On considère la suite $(u_n)$ définie par $u_0= \ln\left (9\right )$, et pour tout entier naturel $n$:

\[u_{n+1}= f\left (u_n\right ).\]

\begin{enumerate}
	\item La fonction $f$ est dérivable sur $\R$ et on a :

	\[\boxed{f'(x) = \dfrac{\textstyle\frac{1}{2}\text{e}^{\frac{x}{2}}}{\text{e}^{\frac{x}{2}}+2}}.\]

	\smallskip

On sait que, pour tout $x \in \R$ on a $\text{e}^{\frac{x}{2}} > 0$, on on déduit que 
$\dfrac12 \text{e}^{\frac{x}{2}} > 0$ et  $\text{e}^{\frac{x}{2}} + 2 > 0$.

Les deux termes du quotient sont supérieurs à zéro, donc  $f'(x)$ est supérieur à zéro sur $\R$, donc la fonction $f$ est strictement croissante sur $\R$.

	\item On sait que si $u > 0$, \: alors $\ln \left(\text{e}^{u}\right) = u$, donc en particulier 
	
\[f\left(2\ln(2)\right) = \ln\left(\text{e}^{\frac{2\ln(2)}{2}} + 2\right) = \ln\left(\text{e}^{\ln(2)} + 2\right) = \ln(2 + 2) = \ln(4) = \ln\left(2^{2}\right) = \boxed{2\ln(2)}.\]

	\item \quad $u_{1} = f\left(u_{0}\right) = \ln\left(\text{e}^{\frac{\ln(9)}{2}} + 2\right) = \ln\left(\text{e}^{\ln\left(9^{\frac{1}{2}}\right)} + 2\right) = \ln\left(\text{e}^{\ln(3)} + 2\right) = \ln(3 + 2) = \boxed{\ln(5)}$
	
	\item Montrons par récurrence que pour tout entier naturel $n$ on a : $2\ln(2) \leqslant u_{n+1} \leqslant u_{n}$.

\begin{list}{\textbullet}{}
\item \textbf{Initialisation}

La fonction $\ln$ est strictement croissante sur $\R$, donc $\ln(4) \leqslant \ln(5) \leqslant \ln(9)$ c'est-à-dire $2\ln(2) \leqslant u_{1} \leqslant u_{0}$. La propriété est donc vérifiée pour $n = 0$.

\item	\textbf{Hérédité}

Supposons que pour $n \in \N$ on ait $2\ln(2) \leqslant u_{n+1} \leqslant u_{n}$.

On sait que la fonction $f$ est strictement croissante donc :

\begin{center}
$2\ln(2) \leqslant u_{n+1} \leqslant u_{n} \iff f\left(2\ln(2)\right) \leqslant f\left(u_{n+1}\right) \leqslant f\left(u_{n}\right) \iff 2\ln(2) \leqslant u_{n+2} \leqslant u_{n+1}$
\end{center}

	Le propriété est donc bien héréditaire.

\item \textbf{Conclusion} 

La propriété est vérifiée pour $n = 0$, et  elle est héréditaire pour tout $n \in \N$  donc d'après le principe de récurrence, on a, pour tout $n \in \N$: $2\ln(2) \leqslant u_{n+1} \leqslant u_{n}$.
\end{list}

	\item D'après la question précédente on a, pour tout $n \in \N$, $u_{n+1} \leqslant u_{n}$, ce qui montre que la suite $\left(u_{n}\right)$ est décroissante.

De plus, pour tout $n \in \N$ on a $u_{n} \geqslant 2\ln(2)$, donc la suite $\left(u_{n}\right)$ est minorée.
	
Donc, d'après le théorème de convergence monotone, la suite $\left(u_{n}\right)$ est convergente vers une limite supérieure ou égale à $2\ln(2)$.
	\item
	\begin{enumerate}
		\item On résout dans $\R$ l'équation $X^{2} - X - 2 = 0$

Calcul du discriminant :

\[\Delta = (-1)^{2} - 4 \times 1 \times (-2) = 1 + 8 = \boxed{9}.\]
$\Delta > 0$ donc il y a deux racines distinctes dans $\R$ :

\[X_{1} = \dfrac{-(-1) - \sqrt{9}}{2 \times 1} = \dfrac{1-3}{2} = -1 \qquad ; \qquad X_{2} = \dfrac{-(-1) + \sqrt{9}}{2 \times 1} = \dfrac{1+3}{2} = 2.\]

		\[\boxed{\mathcal{S} = \left\{-1~;~2 \strut\right\}}\]

		\item \quad $\text{e}^{x} - \text{e}^{\frac{x}{2}} - 2 = 0 \iff \text{e}^{\frac{2x}{2}} - \text{e}^{\frac{x}{2}} - 2 = 0 \iff \left(\text{e}^{\frac{x}{2}}\right)^{2} - \text{e}^{\frac{x}{2}} - 2 = 0$

Posons $X = \text{e}^{\frac{x}{2}}$. On a alors :

\[\left(\text{e}^{\frac{x}{2}}\right)^{2} - \text{e}^{\frac{x}{2}} - 2 = 0 \iff X^{2} - X - 2 = 0\]

D'après la question précédente, les solutions de l'équation $X^{2} - X - 2 = 0$ sont

$X_{1} = -1$ et $X_{2} = 2$. On a donc :

\begin{list}{\textbullet}{}
\item $\text{e}^{\frac{x_{1}}{2}} = -1 \quad (\text{impossible car} \:\text{e}^{\frac{x}{2}} > 0)$ 
\item $\text{e}^{\frac{x_{2}}{2}} = 2 \iff \dfrac{x_{2}}{2} = \ln(2) \iff x_{2} = 2\ln(2)$
\end{list}

\[\boxed{\mathcal{S} = \left\{2\ln (2)\strut\right\}}\]

		\item On a vu à la question \textbf{5.} que la suite $\left(u_{n}\right)$ converge. Soit $\ell$ la limite de la suite $\left(u_{n}\right)$. Comme $f$ est une fonction continue, d'après le théorème du point fixe, on a $\ell$ qui vérifie $f(\ell) = \ell$.

On considère donc l'équation $f(x) = x$ :

\[f(x) = x \iff \ln\left(\text{e}^{\frac{x}{2}} + 2\right) = x \iff \text{e}^{\frac{x}{2}} + 2 = \text{e}^{x} \iff \text{e}^{x} - \text{e}^{\frac{x}{2}} - 2 = 0.\]

D'après la question précédente la solution de cette équation est $x = 2\ln(2)$, donc :

\[\boxed{\lim_{n \to +\infty}u_{n} = 2\ln(2)}.\]
	\end{enumerate}
\end{enumerate}


\bigskip

\end{document}
